% RESEARCH_DRAFT_SHA256: 9187483d786f32c03b8a3258a8ff026640a591493e317902f69103ad0faac0af
% Layout and source-reference wrappers expanded for the web renderer.
\documentclass[11pt]{article}
\usepackage[margin=25mm]{geometry}
\usepackage{fontspec,unicode-math}
\setmainfont{Latin Modern Roman}
\setmathfont{Latin Modern Math}
\usepackage{amsmath,amssymb,mathtools,xurl,hyperref}
\hypersetup{colorlinks=true,urlcolor=blue,linkcolor=blue}
\setlength{\parindent}{0pt}
\setlength{\parskip}{5pt}
\setlength{\emergencystretch}{3em}
\newcommand{\sref}[2]{\hyperlink{source:#1:#2}{[S#2]}}
\newcommand{\cataloguescope}[1]{\paragraph{Source question.}#1}
\newcommand{\floor}[1]{\left\lfloor#1\right\rfloor}
\begin{document}

\section*{Probabilistic boxes: a source-definition audit with three distinct memory models}
\noindent catalogue entry; catalogue entry. Research attempt dated 11 October 2026.
\paragraph{Result in scope.}This manuscript contains a source-specific attempt and complete proofs of the stated auxiliary calculations. It makes no new solution or novelty claim for the entire target.
\section{Original problem and preserved definition}
The following definition and source list are reproduced from the repository. Any earlier inspection dates and reported literature claims in this preserved text describe that source record, not additional verification by this attempt.
% BEGIN PRESERVED DEFINITION
\subsection{Definitions and mathematical statement}

\paragraph{Shared notation.}Standard number systems are $\mathbb N,\mathbb Z,\mathbb Q,\mathbb R,\mathbb C$. Specialized notation and index ranges follow the cited source and the definitions below.


A box model requires a specified joint or conditional probability law for the variables it represents. A probabilistic extension replaces fixed outcomes by random variables.
\paragraph{Record scope and qualifications.}
The catalogue title is retained as metadata. Named certificates, contracts, bounds or frameworks in that title are not asserted as results of the original source extract.
The supplied registry flags this extraction as unrecoverable. It is displayed as a literal source fragment; the definitions provide context rather than a replacement formal statement.

The boxes, variables, observation model, and mathematical target are absent.

Source-record limitation: Source flags the statement as unrecoverable: The linked report was inspected: this is a rhetorical transition immediately followed by a definition of quantum functions, not a freestanding open problem.
The source text below is rendered literally where necessary; damaged formulas and composite extracts are not silently repaired into a new mathematical claim.
The supplied catalogue marks this record as unrecoverable. Its extract and terminology are documented, but a complete original problem has not been reconstructed.
\paragraph{Statement recorded in the supplied source.}
What if we consider probabilistic boxes? The report passage is retained verbatim as the relevant context for this extracted problem. \href{https://doi.org/10.4171/owr/2021/54}{[S2]}
\paragraph{Source status.}
The supplied catalogue labels this source record open; that metadata was checked against the supplied snapshot on 2026-10-11. The label is preserved for provenance and is not an independent verification of present open status or mathematical correctness. The catalogue literature assessment is dated 2026-08-22.
\subsection{Short English statement}
The fragment asks about a probabilistic extension of an unidentified box model.
\subsection{Sources}
\begin{itemize}
\item[S1] Ulam AI and the UnsolvedMath Contributors, UnsolvedMath: A Curated Collection of Open Mathematics Problems, catalogue entry. Catalogue curation and metadata: CC BY 4.0. The source extract is retained; definitions, English summaries and qualifications are editorial additions.
\url{https://huggingface.co/datasets/ulamai/UnsolvedMath}.
\item[S2] Complexity Theory (2022), Oberwolfach Reports, DOI 10.4171/owr/2021/54, pp. 2955--3005
\url{https://doi.org/10.4171/owr/2021/54}.
\item Source-specific Oberwolfach Reports licences, including CC BY-SA 4.0 where applicable, are retained; see THIRD\_PARTY\_NOTICES.md.
\end{itemize}
% END PRESERVED DEFINITION
\subsection{Verbatim supplied catalogue question}
What if we consider probabilistic boxes? The report passage is retained verbatim as the relevant context for this extracted problem.

\section{The fragment in its original setting}
The preserved extract is a rhetorical transition on printed page
2972, immediately before the definition of quantum functions.
It does not ask for a theorem with specified quantifiers. The
primary text supplies finite-dimensional positive operators
$M_{a\mid x}$ with $\sum_aM_{a\mid x}^2=I$ and sequential
probabilities
\[
 p(a_1,\ldots,a_q\mid x_1,\ldots,x_q)
 =d^{-1}\|M_{a_1\mid x_1}\cdots M_{a_q\mid x_q}\|_{\rm F}^2.
\]
Repeated queries are allowed. We investigate the proposed move
from arbitrary probabilistic boxes to this model and identify
what it does and does not represent. This is a source-specific
mathematical audit, not a reconstructed open problem.

\section{Validity of the probability model}
Every term is nonnegative. Summing over the last outcome gives
\[
 \sum_{a_q}\|CM_{a_q\mid x_q}\|_{\rm F}^2
 =\operatorname{tr}\left(C\left(\sum_{a_q}M_{a_q\mid x_q}^2\right)C^*\right)
 =\|C\|_{\rm F}^2,
\]
where $C$ is the preceding product. Iterating and dividing by
$d$ gives total probability one and consistent prefix marginals.
The same calculation works for adaptive next queries because
normalization holds separately for every chosen input.

\section{What commuting positive measurements actually describe}
For finite input and output sets, a commuting family of Hermitian
matrices has a common orthonormal eigenbasis. One can prove this
by diagonalizing one matrix and noting that its eigenspaces are
invariant under all matrices commuting with it; repeatedly
diagonalize restrictions, terminating after finitely many splits.
Write the nonnegative diagonal entries as $m_{a\mid x}(\ell)$,
where $1\le\ell\le d$. The normalization becomes
$\sum_a m_{a\mid x}(\ell)^2=1$, and the sequential law is
\[
 p(a_1,\ldots,a_q\mid x_1,\ldots,x_q)
 =\frac1d\sum_{\ell=1}^d\prod_{j=1}^q m_{a_j\mid x_j}(\ell)^2.
\]
Thus one first samples a uniform latent coordinate $\ell$, then
uses fresh conditionally independent draws with probabilities
$m_{a\mid x}(\ell)^2$ at each query. For adaptive queries the
same conditional formula follows by multiplying the appropriate
probabilities along the branch. The latent coordinate remains
fixed, but a repeated input need not yield a fixed output.

This matters because the report's subsequent Theorem1.1 describes
a mixture of fixed deterministic functions. Under the literal
positive-operator definition this conclusion has a counterexample:
take one input, two outputs, $d=1$, and
\[
 M_0=M_1=1/\sqrt2.
\]
They commute and satisfy the displayed normalization. Every
length-$q$ bit sequence has probability $2^{-q}$; in particular
two repetitions disagree with probability $1/2$. A mixture of
fixed functions on a one-point domain always agrees on repetitions.
The mismatch is in the printed definition/convention, not a
counterexample to a theorem under an unstated projectivity promise.
The root Codex GPT-6 model also read the original definition and
repeated-query convention and checked this counterexample; this
is model-assisted written review, not human verification.

The projective restriction repairs this particular conclusion.
Commuting projectors have common diagonal entries in $\{0,1\}$,
with exactly one output equal to one for every input and coordinate.
Each coordinate is therefore a deterministic function, and the
formula above is their uniform mixture. Conversely, in the
commuting positive model, perfect repeatability for every input
forces this point-mass situation: the repeat-agreement probability
for input $x$ is
\[
 \frac1d\sum_\ell\sum_a p_a(\ell)^2,
 \qquad p_a(\ell)=m_{a\mid x}(\ell)^2.
\]
Each inner sum is at most one, and equality holds only when one
$p_a(\ell)$ equals one, since
$1-\sum_a p_a^2=\sum_{a\ne b}p_ap_b$. An average equal to one
therefore forces every coordinate to be a point mass.

\section{Neither one-query marginals nor arbitrary memory suffice}
A second valid quantum box has $d=2$ and
$M_0=\operatorname{diag}(1,0)$,
$M_1=\operatorname{diag}(0,1)$. Its first output is a fair bit,
as for the scalar box, but it repeats that bit forever. Its
length-$q$ law is supported on the two constant sequences,
each of probability $1/2$. Thus identical one-query marginals
do not identify the memory model or its higher moments.

Conversely, an arbitrary classical history-dependent box can
produce $0$ on its first query and $1$ on its second query to
the same input. No box in the literal finite tracial positive
model realizes this deterministic alternation. Probability zero
for output $1$ on the first query means
$d^{-1}\operatorname{tr}(M_1^2)=0$. Positivity makes all its
eigenvalues zero, so $M_1=0$, and every later outcome sequence
containing $1$ has probability zero as well. Additional
instruments, internal updates or an alternative initial-state
model would change this restriction.

\section{Attempted inference, audit, and status}
The attempted identification of all probabilistic boxes, commuting
quantum functions, and mixtures of fixed functions fails at the
explicit examples above. The audit verified the source's use of
positive operators rather than projectors and of repeated query
sequences, and checked normalization and adaptivity directly.
No computation or theorem prover was used. The source fragment
still lacks a freestanding unresolved mathematical target; these
results document its precise context without inventing one.

\section{Process and contributions}
Mathematics and model-assisted adversarial self-audit: OpenAI gpt-6-astra, ultra reasoning, as dispatched for this research batch; exact serving revision not exposed. Codex (GPT-6; exact serving revision not exposed) materially assisted with orchestration and format review. No human mathematical authorship, mathematical edits, or independent human proof checking is asserted.
The coordinating Codex (GPT-6; exact serving revision not exposed) also provided a separate model-assisted written-proof review of this manuscript. This is not independent human verification.
Contribution prompt A was applied to the full definition and supplied source row, then prompt B was applied adversarially to the derived results. The item-specific audit above records the scope checks and rejected extrapolations. Local assembly and structural checks do not establish mathematical correctness. The source catalogue status is preserved in TOP\_PROBLEM and is separate from this attempt status.
\paragraph{Integration review.} Codex (GPT-6; exact serving revision not exposed) reviewed the written arguments and expanded layout and source-reference wrappers. Repository metadata and the production web renderer were checked using Python 3.13 and JavaScript with MathJax 3.2.2. This model-assisted review is not independent human proof verification; no PDF compilation or formal proof checking is asserted.

\section{Completion Estimate}
\noindent\textbf{COMPLETION: 0\%}
The proved material is a rederivation of known or auxiliary facts rather than a substantiated new advance on the original target. The preserved source fragment is a rhetorical transition, not a quantified open problem. Three explicit box models and the literal-definition mismatch do not create a replacement target.
\section{References}
The complete original Sources subsection above is retained, including catalogue attribution and licence notices. The additional list records primary sources actually consulted for this attempt on 11 October 2026; uncited source links reproduced above were not all rechecked.
\begin{itemize}
\item Henry Yuen, Noncommutative Property Testing, OWR54/2021 pp.2971–2974; full primary source passage read in cached text, and page2972 operator definition visually checked from the primary PDF. The original agenda, stated known results and convention mismatch are distinguished. \url{https://ems.press/content/serial-article-files/46932}.
\end{itemize}
\section{Final status}
\textbf{Unresolved in this attempt.} The preserved source fragment is a rhetorical transition, not a quantified open problem. Three explicit box models and the literal-definition mismatch do not create a replacement target.
\end{document}
